% !TeX root = ../main.tex
\begin{frame}{课程评价——吴莹钧}
  \begin{SlideCanvas}
    \node[anchor=north west,inner sep=0pt,text width=850pt,align=left,
      text=BICNavy] at (55,80) {\begin{minipage}{850pt}
      \usebeamerfont{normal text}\raggedright
      \setlength{\parindent}{0pt}\setlength{\parskip}{4pt}
      \setlength{\columnsep}{26pt}\setlength{\multicolsep}{0pt}
      \clubpenalty=10000\widowpenalty=10000
\par\addvspace{6pt}{\usebeamerfont{BIC point}一、课程评价\par}\nobreak
该课程涵盖医学影像、智能手术机器人、柔性电子、MEMS等多个前沿方向，内容丰富，具有较强的交叉性和应用性。授课教师结合自身科研工作，介绍了相关技术原理和实际应用，让我对医工交叉领域的研究现状和发展趋势有了更直观的了解。\par
\par\addvspace{6pt}{\usebeamerfont{BIC point}二、学习收获\par}\nobreak
我印象比较深刻的是磁共振成像、全数字PET及柔性电子等内容。通过学习，我认识到医工交叉不仅是将工程技术应用于医学，还需要从实际临床需求出发，综合考虑器件性能、系统集成和应用条件。对于集成电路背景的学生来说，这门课程也让我对自己的研究方向有了更多新的思考。\par
\par\addvspace{6pt}{\usebeamerfont{BIC point}三、意见与建议\par}\nobreak
我认为课程整体安排比较合理，但部分专题专业性较强，对于不同学科背景的学生来说，理解起来有一定难度。建议授课时适当补充基础知识，并结合具体案例讲解临床需求与工程技术之间的联系。此外，也可以增加一些课堂交流和讨论，让学生从不同学科的角度思考问题，更好地体现医工交叉的特点。\par
    \end{minipage}};
  \end{SlideCanvas}
  \note{原始材料为用户提供的四份课程反馈之一，正文逐字保留，仅调整标题、段落和编号排版。吴莹钧《医工交叉研究前沿进展》课程反馈与建议。一、课程评价\par 该课程涵盖医学影像、智能手术机器人、柔性电子、MEMS等多个前沿方向，内容丰富，具有较强的交叉性和应用性。授课教师结合自身科研工作，介绍了相关技术原理和实际应用，让我对医工交叉领域的研究现状和发展趋势有了更直观的了解。\par \par 二、学习收获\par 我印象比较深刻的是磁共振成像、全数字PET及柔性电子等内容。通过学习，我认识到医工交叉不仅是将工程技术应用于医学，还需要从实际临床需求出发，综合考虑器件性能、系统集成和应用条件。对于集成电路背景的学生来说，这门课程也让我对自己的研究方向有了更多新的思考。\par \par 三、意见与建议\par 我认为课程整体安排比较合理，但部分专题专业性较强，对于不同学科背景的学生来说，理解起来有一定难度。建议授课时适当补充基础知识，并结合具体案例讲解临床需求与工程技术之间的联系。此外，也可以增加一些课堂交流和讨论，让学生从不同学科的角度思考问题，更好地体现医工交叉的特点。}
\end{frame}
